\BourbakiBackSection{历史注记}

（括号中的数字指本注记末尾的参考文献。）

极限与连续的观念可以追溯到古代；若不从这个观点出发，不仅对希腊数学家、而且对希腊哲学家、特别是亚里士多德进行系统的研究，就无法写出一部完整的极限与连续观念史。此外，还必须追溯这些观念在文艺复兴时期数学中以及 在微分学与积分学开端时期的演变。这样一项研究，尽管进行起来无疑会引人入胜，却将远远超出本注记的范围。

应当认为是 Riemann 创立了拓扑学，正如现代数学的许多其他分支一样。他第一个尝试表述拓扑空间的概念；他孕育了关于这种空间的自足理论的设想；他定义了在拓扑学后来的发展中起卓越作用的不变量（“Betti 数”）；他第一个把拓扑学应用于分析（Abel 积分的周期）。但是，十九世纪上半叶的思想潮流已经从不止一个方面为 Riemann 铺平了道路。首先，把数学建立在坚实基础之上的愿望——它是整个十九世纪直至今天许多重要研究的缘由——使人们正确理解了收敛级数与趋于极限的数列的概念（Cauchy、Abel），以及连续函数的概念（Bolzano、Cauchy）。另一方面，由 Gauss 与 Argand 提出的复数（或者像以前人们所称的那样，“虚”数甚至“不可能”数）的几何表示（以平面上的点表示），已为大多数数学家所熟悉；这一进步与本世纪在 Hilbert 空间研究中采用几何语言具有同一数量级，并且蕴含着对每个能连续变化的对象进行几何表示这一可能性的萌芽。至于 Gauss，他本来就会由于其关于几何基础、非欧几何与曲面的研究而自然地达到这类概念，而且他似乎早已想到这种可能性，因为他在（独立于 Argand 与法国数学家）定义复数的几何表示时使用了“二次延展的量”一语（{[}1{]}，第 101-103 页及第 175-178 页）。

Riemann 关于代数函数及其积分的工作，以及他对几何基础的思考（在很大程度上来自他对 Gauss 著作的研究），使他制定了一个研究纲领，这恰好就是现代拓扑学的纲领，并使他开始实现这一纲领。例如，他在其 Abel 函数论中说道（{[}2{]}，第 91 页）：

“在研究由恰当微分的积分所得到的函数时，Analysis situs 的某些定理几乎是不可或缺的。这个由 Leibnitz 用过的名称（尽管也许是在多少有些不同的意义上）应当称呼连续量理论中这样的部分：它研究这些量时，不是脱离它们的位置、以它们彼此度量来研究，而是撇开一切度量观念，只考虑它们的位置关系与包含关系。我把它完全脱离一切度量的研究留给以后的机会\ldots{}”

而在其著名的就职演说“论几何学所赖以成立的假设”（{[}2{]}，第 272 页）中：

“\ldots{} 多次延展的量的一般概念（(*)）——它把空间量概念作为特例包含在内——至今完全没有得到研究\ldots”（第 272 页）

“\ldots{} 量的概念以一个元素可作各种不同的规定为前提。按照能否经由连续的过渡过程从一个规定过渡到另一个规定，这些规定构成一个连续流形或一个离散流形：在前一情形，这些规定称为该流形的点\ldots”（第 273 页）

“\ldots{} 度量在于把待比较的量叠合起来，因此为了度量，我们需要某种以一个量作为另一个量的尺度的手段。缺乏这种手段时，只有当一个量是另一个量的一部分时，我们才能比较这两个量\ldots{} 在此情形下所能进行的研究构成量理论中独立于度量理论的一部分，其中诸量不被看作独立于其位置而存在，也不被看作可用度量单位表出，而是被看作流形中的区域。这样的研究在数学的若干部分中已成为必要，特别是在多值解析函数论中\ldots”（第 274 页）

“\ldots{} 在给定的流形中，只要这是可能的，位置的确定都可归结为有限多个数值规定。然而存在这样的流形，其中位置的确定需要的不是有限多个规定，而是一个无穷序列、甚至是量规定的一个连续流形。例如，一个函数在给定区域上的各种可能规定，或者一个空间图形的各种可能形状，都给出这种类型的流形。”（第 276 页）

请注意最后这段话中函数空间研究的最初观念；Riemann 在其博士论文中已经表述过同一观念：他在谈到被称为 Dirichlet 原理的极小问题时说道，“这些函数的全体”“构成一个自身封闭的连通区域”（{[}2{]}，第 30 页）；这一表述虽不完善，却正是 Hilbert 后来给出的 Dirichlet 原理的证明、以及函数空间对变分法的大部分应用的萌芽。

如前所述，Riemann 开始实施这一宏伟纲领：他定义了“Betti 数”，先是对曲面（{[}2{]}，第 92-93 页），后来（{[}2{]}，第 479-482 页；另参看{[}3{]}）对任意维数的流形，并把这一定义应用于积分论；关于这一点以及这一理论自 Riemann 时代以来的巨大发展，请读者参看本丛书中代数拓扑各章的历史注记。

在 Riemann 所设想的那种拓扑空间的一般理论得以发展之前，实数理论、数集理论以及直线上、平面上与空间中的点集理论，都必须比 Riemann 时代得到更为系统的研究。这些研究另一方面与关于无理数本性的探讨（Bolzano 的半哲学式探讨与 Dedekind 的本质上数学的探讨）以及实变函数论中的进展相联系；在实变函数论中，Riemann 本人以其积分定义与其三角级数理论作出了重要贡献，du Bois-Reymond、Dini 与 Weierstrass 等人也对它有所贡献；这些工作是十九世纪下半叶的工作，尤其是 Cantor 的工作，他第一个（最初在直线上，后来在 n 维欧氏空间中）定义了聚点、闭集、开集、完备集诸概念，并得到了关于这些集合在直线上的结构的主要结果（参看第四章的历史注记）。在这方面，不仅应参看 Cantor 的著作{[}4{]}，还应参看他与 Dedekind 的极其有趣的通信{[}5{]}，其中可以找到维数作为拓扑不变量这一思想的明确表述。该理论后来的进展在 Schoenflies 的书{[}6{]}中以半历史、半系统的方式加以追溯；迄今最重要的收获是 Borel-Lebesgue 定理，即欧氏 \(n\) 维空间 \(\mathbf{R}^n\)（参看第六章，§ 1）的每个有界闭子集都满足本章 § 9 的公理（C \(^{''}\)）这一事实（该定理最初是 Borel 对直线上的一个闭区间以及覆盖它的可数开区间族证明的）。

Cantor 的思想最初曾遭到激烈的反对（参看第一卷第一章至第四章的历史注记）。无论如何，他关于直线上与平面上的点集的理论很快就被法国与德国的函数论学派（Jordan、Poincaré、Klein、Mittag-Leffler，以及后来的 Hadamard、Borel、Baire、Lebesgue 等）所利用和传播；特别是 Borel 早期的每部教程都含有这一理论的初等阐述（例如见{[}7{]}）。随着这些思想的传播，把它们应用于不是由点组成、而是由曲线或函数组成的集合的可能性，开始在各个方面被考虑，其见证是 Ascoli 1883 年一篇论文的标题“论一簇曲线的极限曲线”{[}8{]}，以及 Hadamard 1896 年在苏黎世数学家大会上的报告{[}9{]}；所有这些都与 Volterra 1887 年引入“线函数”以及创立“泛函演算”即自变量为函数的函数论密切相关（参看 Volterra 关于泛函分析的书{[}10{]}）。另一方面，Hilbert 的著名论文{[}11{]}——他在其中重新采纳 Riemann 的思想，证明了 Dirichlet 原理中最小值的存在，并开创了变分法中的“直接方法”——清楚地表明了考虑使 Bolzano-Weierstrass 原理成立、也就是说使每个序列都有收敛子列的函数集的重要性。这类集合无论如何已开始起重要作用，不仅在变分法中，而且在实变函数论中（Ascoli、Arzelà），稍后还在复变函数论中（Vitali、Carathéodory、Montel）。最后，函数方程的研究、特别是 Fredholm 对冠以其名的那类方程的求解，使把函数看作变元、把函数集看作点集成为常见之事，并且使在这一场合使用几何语言如同在 \(n\) 维欧氏空间中一样自然（这种空间同样逃避“直观”，因此长期是许多数学家不信任的对象）。特别是 Hilbert 关于积分方程的值得纪念的工作{[}12{]}导致 Erhard Schmidt{[}14{]}对 Hilbert 空间的定义与几何研究，与欧氏几何完全类似。

与此同时，公理理论的概念由于几何基础方面的大量工作而变得越来越重要；在这方面，Hilbert 的贡献{[}13{]}具有特别决定性的影响。在这项工作过程中，Hilbert 于 1902 年（{[}13{]}，第 180 页）给出了 Riemann 意义下“二次延展流形”的第一个公理化定义，Hilbert 说这个定义构成“对 Analysis situs 作严格公理处理的基石”。Hilbert 还使用了邻域（其意义受他把自己局限于其中的问题的要求的限制）。

把点集性质与函数集性质的共同内容抽象出来的最初尝试应归功于 Fréchet{[}15{]}与 F. Riesz{[}16{]}。前者从可数极限的概念出发，未能构造出一个方便而富有成效的公理系统，但至少他认识到了 Bolzano-Weierstrass 原理{[}即 § 9 的公理（C）限于可数序列的情形{]}与 Borel-Lebesgue 定理{[}§ 9 的公理（C''）{]}之间的关系；正是在这一场合他引入了“紧”这个词，尽管其意义与本丛书所用者略有不同。至于 F. Riesz，他以聚点概念（或者不如说“导集”概念，这二者是一回事）为出发点，他的理论仍不完全，且只以概要的形式出现。

今天所理解的一般拓扑学始于 Hausdorff（{[}17{]}，第 7、8、9 章），他重新采用了邻域概念（他所谓邻域，即本丛书术语中所称的“开邻域”），并从 Hilbert 关于平面上邻域的公理中挑选出那些能使其理论具有所期望的全部精确性与一般性的公理。他作为出发点的公理本质上是（考虑到他的邻域概念与我们的邻域概念之间的差别）§ I 的公理（ \(V_{I}\)）、（ \(V_{II}\)）、（ \(V_{III}\)）、（ \(V_{IV}\)）以及 § 8 的（H），而他推演这些公理的推论的那一章一直是公理理论的一个典范：抽象而又预先适应于应用。Hausdorff 的工作自然是后来一般拓扑学研究、特别是莫斯科学派工作的出发点，后者的工作大部分指向度量化问题（参看第九章的历史注记）；这里我们特别回顾 Alexandroff 与 Urysohn 对紧空间的定义（以“双紧空间”之名），以及 Tychonoff 关于紧空间的积仍为紧空间的证明{[}19{]}。最后，H. Cartan{[}20{]}引入滤子给拓扑学带来了一件宝贵的工具，可用于各种各样的应用（在其中它比“Moore-Smith 收敛”概念{[}18{]}更为有利）。此外，超滤子定理（定理 1，§ 6）的发展阐明并简化了理论。

